\section{Messaging and File-Transfer Workflows}
\subsection{Private and broadcast chat}

For a private chat, the server forwards \texttt{CHAT} only to the named recipient; for a broadcast, it uses the online recipient snapshot captured when the message is enqueued. Each outgoing chat has a stable UUID. A hop-level UDP ACK confirms that a DATA frame reached the next transport endpoint, \emph{not} that the recipient application accepted it. Only the recipient's encrypted \texttt{CHAT\_ACK} causes a final \texttt{CHAT\_RECEIPT (DELIVERED)} at the sender. \Cref{fig:chatseq} illustrates one private message.

\begin{figure}[H]
\centering\begin{tikzpicture}[>=Latex,font=\small,
 participant/.style={draw=black!65,rounded corners=2pt,fill=gray!7,minimum width=2.45cm,minimum height=0.6cm,align=center},
 msg/.style={->,thick}, ack/.style={->,dashed,draw=black!75}]
\node[participant] (a) at (0,0) {Alice (sender)};
\node[participant] (s) at (6.4,0) {Collaboration server};
\node[participant] (b) at (12.8,0) {Bob (receiver)};
\foreach \x in {0,6.4,12.8}{\draw[densely dotted,black!55] (\x,-.35)--(\x,-5.85);}
\draw[msg] (0,-.90)--node[above,font=\scriptsize]{CHAT + stable message\_id}(6.4,-.90);
\draw[ack] (6.4,-1.55)--node[above,font=\scriptsize]{UDP transport ACK}(0,-1.55);
\draw[msg] (6.4,-2.20)--node[above,font=\scriptsize]{Forward encrypted CHAT}(12.8,-2.20);
\draw[ack] (12.8,-2.85)--node[above,font=\scriptsize]{UDP transport ACK}(6.4,-2.85);
\node[align=center,fill=white,font=\scriptsize] at (12.8,-3.52) {De-duplicate\\and accept};
\draw[msg] (12.8,-4.18)--node[above,font=\scriptsize]{CHAT\_ACK (application)}(6.4,-4.18);
\draw[ack] (6.4,-4.85)--node[above,font=\scriptsize]{CHAT\_RECEIPT: DELIVERED}(0,-4.85);
\node[font=\scriptsize,align=center,text=black!70] at (6.4,-5.52) {Sender marks delivered only after the recipient application acknowledges.};
\end{tikzpicture}

\caption{Chat delivery: transport acknowledgments and application acceptance are distinct.}
\label{fig:chatseq}
\end{figure}

\subsection{Reliable UDP file transfer}

\texttt{FileSender} calculates SHA-256 and emits \texttt{FILE\_START} with filename, size, chunk count and checksum. It reads the file in 800-byte blocks rather than loading it into memory; the server relays each encrypted chunk. The receiver checks chunk index and total size, writes a temporary \texttt{.part} file, then verifies SHA-256 before making the completed file visible. The sender receives a final \texttt{FILE\_RECEIPT}. Interrupted transfers remove incomplete temporary files, while verified downloads are retained.

For example, a 32-KiB file contains \(\lceil 32768/800\rceil=41\) data chunks and two boundary messages (\texttt{FILE\_START}, \texttt{FILE\_END}): \textbf{43 application DATA packets} before retransmissions. Each hop has its own UDP ACK exchange.

\begin{figure}[H]
\centering\begin{tikzpicture}[>=Latex,font=\small,
 participant/.style={draw=black!65,rounded corners=2pt,fill=gray!7,minimum width=2.6cm,minimum height=0.6cm,align=center},
 msg/.style={->,thick}, ack/.style={->,dashed,draw=black!75}]
\node[participant] (a) at (0,0) {Sender};
\node[participant] (s) at (6.4,0) {Server relay};
\node[participant] (b) at (12.8,0) {Recipient};
\foreach \x in {0,6.4,12.8}{\draw[densely dotted,black!55] (\x,-.35)--(\x,-6.8);}
\draw[msg] (0,-.85)--node[above,font=\scriptsize]{FILE\_START (name, size, SHA-256)}(6.4,-.85);
\draw[msg] (6.4,-1.42)--node[above,font=\scriptsize]{Relay FILE\_START}(12.8,-1.42);
\draw[msg] (0,-2.16)--node[above,font=\scriptsize]{FILE\_CHUNK (index, up to 800 B)}(6.4,-2.16);
\draw[ack] (6.4,-2.68)--node[above,font=\scriptsize]{UDP ACK}(0,-2.68);
\draw[msg] (6.4,-3.22)--node[above,font=\scriptsize]{Relay chunk on file stream}(12.8,-3.22);
\draw[ack] (12.8,-3.74)--node[above,font=\scriptsize]{UDP ACK}(6.4,-3.74);
\node[font=\scriptsize,text=blue!65!black] at (10.3,-4.22) {repeat chunks};
\draw[msg] (0,-4.75)--node[above,font=\scriptsize]{FILE\_END}(6.4,-4.75);
\draw[msg] (6.4,-5.28)--node[above,font=\scriptsize]{Relay FILE\_END}(12.8,-5.28);
\node[font=\scriptsize,align=center,fill=white] at (12.8,-5.77) {Verify size\\and SHA-256};
\draw[msg] (12.8,-6.30)--node[above,font=\scriptsize]{FILE\_RECEIPT (ok / error)}(6.4,-6.30);
\draw[msg] (6.4,-6.73)--node[above,font=\scriptsize]{Forward FILE\_RECEIPT}(0,-6.73);
\end{tikzpicture}

\caption{File lifecycle with chunk relay, per-hop ACKs and end-to-end hash receipt.}
\label{fig:fileseq}
\end{figure}

\subsection{Disconnection and pending-message recovery}

A chat whose delivery is uncertain remains in a per-account RAM outbox with status \texttt{PENDING} or \texttt{SENDING}. After a server restart, demo mode reconnects automatically; password account mode asks the user to authenticate again. The sender then replays pending UUIDs in order, and the receiver suppresses duplicates by sender and message ID. A deliberate logout clears the displayed chat history and terminates the session, but the same account can resume its pending outbox \emph{while the local client process remains alive}. Closing that process loses the RAM-only queue; offline message composition is not implemented.

\begin{figure}[H]
\centering
\evidenceslot{03_pending_before_restart.png}{Two pending messages}{Controlled CHAT DATA loss while Alice was connected.}{12cm}
\caption{Actual account-mode sender before server restart: two queued messages remain PENDING and delivery is uncertain. Both were composed while connected.}
\label{fig:pending-evidence}
\end{figure}

\begin{figure}[H]
\centering
\evidenceslot{04_delivered_after_relogin.png}{Ordered replay after relogin}{Alice has no pending items; Bob receives each original UUID once.}{13cm}
\caption{Actual recovery after restarting the normal server and reauthenticating Bob, then Alice. Both local UI processes stayed alive. Logs confirm the original UUIDs, order A then B, one receipt per message and an empty sender outbox.}
\label{fig:recovery-evidence}
\end{figure}
